\chapter{Literature Review}

\section{Customer Feedback and Aspect-Based Sentiment Analysis}

Customer-review mining seeks to convert unstructured opinions into structured
evidence about product or service features. Early feature-based review
summarisation already separated what was discussed from the polarity expressed
towards it \cite{hu2004mining}, while broader sentiment-analysis work
distinguished document-level judgement from more fine-grained opinion mining
\cite{pang2008opinion}. ABSA develops this distinction by associating a
sentiment with an aspect term or category. Shared tasks such as SemEval-2016
formalised aspect-category and polarity settings across several domains
\cite{pontiki2016semeval}, and later surveys organise ABSA into aspect
extraction, aspect-category detection, and aspect-level sentiment
classification \cite{zhang2023survey}.

This distinction is methodologically important for FABSA. A multi-aspect review
cannot be represented faithfully by attaching one global sentiment to every
detected aspect. Aspect-level architectures therefore condition sentiment on
the target aspect, for example by learning aspect-specific context
representations \cite{tang2016deep}. The present dissertation adopts the same
principle throughout: the atomic claim combines a review, one candidate aspect,
and one candidate sentiment, and review-level predictions are formed by
unioning the applicable claims.

\section{Candidate-Conditioned and Multi-Label ABSA}

Several ABSA formulations support treating the aspect as textual input rather
than as an implicit fixed output index. Sun et al.\ construct an auxiliary
sentence from an aspect and transform ABSA into a sentence-pair problem related
to question answering and natural-language inference
\cite{sun2019auxiliary}. Seoh et al.\ similarly use natural-language prompts to
support arbitrary aspect targets in zero- and few-shot sentiment classification
\cite{seoh2021open}. These studies motivate a review--candidate cross-encoder:
the model can condition its decision on the supplied aspect semantics instead
of learning only a document representation.

Other work represents the output directly in natural language. Unified and
text-generation formulations allow several ABSA subtasks or
aspect--sentiment outputs to share one generative interface
\cite{yan2021unified}. Instruction-based ABSA extends this principle to
prompted task definitions \cite{scaria2024instructabsa}. Such work supports the
candidate-pair Qwen interface used here, but does not imply that free-form
generation is required. A one-token applicability score is deliberately used
in the core experiments to keep thresholds, empty predictions, and
frozen-to-adapted comparisons auditable.

Weak-supervision research demonstrates that label-side semantics can support
ABSA without conventional labelled examples. Huang et al.\ use small sets of
keywords describing aspects and sentiments to learn joint aspect--sentiment
topic representations \cite{huang2020weakly}. AX-MABSA begins from one
indicative word per class and produces multi-label aspect-category--sentiment
sets \cite{kamila2022ax}. These approaches establish that label names or
keywords are genuine supervision. They are therefore adjacent to, but stronger
than, this dissertation's minimal-description treatment: automatically mined
corpus keywords are excluded from the strict description regime.

\section{From Fixed Labels to Evolving Taxonomies}

Conventional supervised text classifiers assume that the output label set is
fixed. Work on evolving label sets and dynamic text classification instead
considers categories added after initial model training
\cite{godbole2005evolvinglabelsets,wohlwend2019dynamictextclassification}.
Open-domain topic-classification systems likewise allow users to provide or
revise a taxonomy \cite{ding2022opendomain}. These settings motivate the
practical problem, but this dissertation makes a narrower claim: the candidate
taxonomy is supplied at inference, and the task is selection from that
canonical set rather than discovery of unrestricted topics.

Yin et al.\ distinguish label-partially-unseen evaluation---training on some
labels and testing on all labels---from label-fully-unseen classification with
no task-specific labelled data \cite{yin2019benchmarking}. This distinction
prevents an important terminology error. The dissertation's Level~2 is
label-partially-unseen because eleven trained aspects and one unseen aspect are
jointly evaluated. Level~1 is an easier singleton-candidate diagnostic, not a
fully-unseen task: the model still learns the task from the other aspects.
Levels~3 and~4 extend the partially unseen setting to two unseen labels and an
unseen parent group.

The joint prediction of seen and unseen labels corresponds to the generalised
zero-shot concern that a model can perform well on familiar labels while
suppressing unfamiliar ones. Generalised zero-shot evaluation conventionally
reports seen and unseen performance separately and uses their harmonic mean to
expose this imbalance \cite{xian2017zsl}. Although that evaluation convention
was standardised in computer vision, the same seen/unseen separation appears
in zero-shot multi-label text classification
\cite{liu2021hierarchy}. Accordingly, Levels~2--4 report overall, seen, unseen,
and harmonic-mean results rather than allowing high seen-label performance to
hide failure on the held-out aspects.

\section{Structured Multi-Label Spaces and Hierarchy}

Multi-label text classification permits several correct labels for one
document, so instance-based and label-based metrics reveal different behaviour
\cite{tsoumakas2007multilabel,zhang2014review}. Zero-shot multi-label research
has additionally exploited known structure over label spaces. Rios and Kavuluru
study few- and zero-shot labels in structured medical taxonomies
\cite{rios2018structured}. Liu et al.\ define both unseen-only and generalised
seen-plus-unseen settings and use parent--child relations to strengthen
multi-label prediction \cite{liu2021hierarchy}. Semantic-knowledge approaches
likewise combine class descriptions and class hierarchy when recognising
unseen categories \cite{zhang2019semantic}.

FABSA's twelve labels contain explicit parent--child structure, for example
\texttt{Staff support: Email} and \texttt{Staff support: Phone}. Keeping the
full hierarchical name visible is therefore not cosmetic: it supplies a
controlled form of label semantics. The Level~4 parent-group holdout uses this
structure as an evaluation intervention. It does not implement hierarchical
inference or assume that a parent prediction entails every child. Instead, it
asks whether models trained without all children of a parent can transfer to
that coherent unseen sub-taxonomy.

\section{Names, Descriptions, and Label-Side Supervision}

Treating labels as text provides an interface for unseen categories. GILE
jointly embeds inputs and textual label descriptions, avoiding a classifier
whose parameters are inseparable from a fixed label inventory
\cite{pappas2019gile}. Meng et al.\ go further in a dataless setting by
starting from label names and using unlabelled documents for language-model
self-training \cite{meng2020labelnames}. These studies justify name-only as a
meaningful semantic baseline, but they also show that information regimes must
be stated: name-only plus target-corpus self-training is not the same treatment
as name-only inference with a frozen encoder.

Label descriptions may add objective definitions, related terms, or natural
language templates. Gao et al.\ curate compact label-description data using
related terms, dictionary definitions, Wikipedia sentences, and
sentiment-specific templates; their experiments show gains over zero-shot
classification and reduced sensitivity to pattern choice
\cite{gao2023labeldesc}. This supports the dissertation's use of a short
taxonomy gloss when no official FABSA annotation guidance is available. Because
the FABSA labels are project-specific hierarchical composites, a uniformly
authored definition is more reproducible than selecting different external web
pages for different labels.

Gao et al.'s intervention is stronger than merely appending a definition at
inference: label terms, related terms, dictionary or Wikipedia text, and
templates become labelled description examples used to fine-tune the model.
The present core protocol does not expose held-out aspects through such
task-specific training. Its secondary rich condition adds fixed aliases and
taxonomy boundaries only at representation time. This preserves the same
training manifest and threshold as the minimal condition, so
\texttt{RR}-minus-\texttt{DD} measures richer label-side guidance rather than
additional target-aspect optimisation.

Description choice itself can become hidden model selection. Basile et al.\
note that a true zero-shot setting lacks a labelled target development set and
therefore rank or aggregate multiple descriptions without target labels
\cite{basile2022ranking}. Chu et al.\ refine weak label representations using
unlabelled target documents and clustering \cite{chu2021labelrefinement}, while
Yano et al.\ generate and rerank keywords for relevance, diversity, and
inter-class exclusivity \cite{yano2024relevance}. Both are useful extensions,
but they consume an additional information source or introduce a richer
description generator. The core thesis instead fixes one minimal definition
using only the canonical names, parent--child hierarchy, and general language
knowledge. Target reviews, corpus statistics, labels, predictions, and error
analysis are prohibited.

The description's usefulness is then an empirical question, not an editing
criterion. Once frozen, name-only and name-plus-definition conditions can be
compared on identical examples. A poor result means that the fixed description
did not help that model under the declared protocol; it does not authorise
rewriting the description after observing target results. The Level~3
\texttt{NN}/\texttt{DN}/\texttt{ND}/\texttt{DD} crossover extends this logic
to incomplete documentation. The two directions are both required because a
one-direction result would confound description availability with the intrinsic
difficulty of the two held-out aspects. The secondary \texttt{RR} condition is
therefore symmetric and deliberately does not expand the experiment to every
possible name/minimal/rich pairing.

\section{A Matched Spectrum of Model Capacity}

The five core methods are chosen to test representational and adaptation
capacity without turning the thesis into an unrestricted model sweep. Strict
train-only TF--IDF provides a lexical baseline whose vocabulary and inverse
document frequencies are fitted inside the training scope. E5-base-v2 provides
a frozen dense semantic baseline; E5 was trained contrastively for transferable
single-vector text representations across retrieval and classification tasks
\cite{wang2022e5}. The supervised DistilBERT cross-encoder supplies contextual
token interaction while retaining a relatively compact encoder. DistilBERT
uses pre-training-stage knowledge distillation to reduce BERT's size and
inference cost \cite{sanh2019distilbert}, while the sentence-pair design follows
the candidate-conditioned ABSA motivation discussed above.

Qwen3 supplies the common causal language-model backbone for the final two
methods \cite{yang2025qwen3}. Frozen Qwen converts the next-token logits for
fixed ``yes'' and ``no'' verbalizers into an applicability score. QLoRA keeps
the same input, verbalizers, score, and base-model revision while adapting
low-rank modules. LoRA freezes the base weights and learns low-rank updates
\cite{hu2022lora}; QLoRA back-propagates through a frozen four-bit model into
those adapters to reduce memory requirements \cite{dettmers2023qlora}.

This matched control is scientifically important. Recent zero-shot multi-label
work reports that supervised fine-tuning can improve overall performance while
eroding unseen-label capability, even when label descriptions are present
\cite{chen2025preserving}. QLoRA is therefore not assumed to win because it is
adapted. The comparison asks whether adaptation improves the candidate-pair
task while preserving the semantic route to aspects excluded from
task-specific training.

\section{Metrics and Statistical Uncertainty}

Multi-label evaluation requires more than one F1 variant. Pair micro-F1
aggregates claim decisions, macro-F1 exposes rare aspect--sentiment labels, and
example-based F1 describes review-level set quality. In all-row unseen-aspect
evaluation, true-negative empty rows can dominate, so presence F1 must be
defined over the positive class rather than counting correct absence as
positive evidence. Exact match, prediction-set size, false-positive and
false-negative rows, and conditional sentiment coverage further separate
aspect detection from sentiment assignment. General metric analyses caution
that different averaging rules answer different questions
\cite{sokolova2009systematic}.

Because every matched system is evaluated on the same reviews and candidate
claims, uncertainty analysis should preserve that pairing. Dror et al.\ show
that statistical procedures in NLP must match the evaluation measure and
experimental structure \cite{dror2018hitchhiker}. Paired evaluation likewise
uses the common instances rather than comparing two unrelated averages
\cite{peyrard2021paired}. The dissertation therefore resamples reviews as
clusters, retains all candidate claims belonging to a sampled review, and
recomputes nonlinear metrics within each replicate. Paired metric-difference
intervals are primary for model and description contrasts. Per-aspect
wins/ties/losses and sign-flip evidence address variation across the registered
taxonomic units, while multiple-comparison claims are kept within
pre-registered families \cite{dror2017replicability}.

\section{Synthesis and Positioning}

The literature provides strong support for the components of the revised
route: ABSA requires aspect-conditioned sentiment; auxiliary-sentence and
prompt methods support candidate-conditioned inputs; zero-shot text
classification distinguishes unseen-only from joint seen/unseen evaluation;
structured multi-label work exploits label descriptions and hierarchy; and
parameter-efficient adaptation supplies a feasible test of task-specific LLM
training.

The dissertation combines these components into one controlled difficulty
ladder for customer-feedback ABSA. Its most distinctive experiment is not a new
description-generation algorithm. It is the matched dual-unseen crossover,
which tests whether incomplete semantic documentation changes competition
between two unseen aspects while the training data, candidate pairs, model, and
budget remain fixed. The parent-group holdout then asks whether the same methods
survive a coherent sub-taxonomy shift. This positioning keeps the contribution
challenging but bounded: supplied candidate labels rather than arbitrary topic
generation, and evidence about degradation, description, and adaptation rather
than a claim of universal zero-shot performance.
